% Review prose, not an edit to the live manuscript. Numerical macros come from
% module 07 and must be regenerated with the accompanying review exhibits.

\section{Collective reparations and local development: results for review}

Collective reparations can improve conditions in communities affected by
violence, but place-based investment need not have the same consequences for
every resident. Better local opportunities may encourage people to remain;
greater resources and connectivity may also make moving more feasible. The
present evidence does not distinguish these competing pathways conclusively.
This review brings the community, household and individual estimates together
without treating related observations as independent replications.

The numerical findings concern the approved geographic sample: Apurimac and
Huancavelica, together with La Convencion province in Cusco and Huancayo
province in Junin. This geography contains 1,162 of the 5,712 communities in
the supplied victim registry. The analyses compare adjacent victimization
categories B and C around the official score threshold of 0.062320. The
geography and common estimation window were retained after historical
exploratory design work; they were not selected prospectively before the
original research. The search history must accompany the eventual paper.

Treatment is cumulative recorded receipt through 2012 for the 2013 outcomes
and through 2016 for the 2017 outcomes. Under the project lead's binding
measurement decision, the CMAN register's project year is assumed to be the
year of both allocation and delivery. This assumption avoids inventing dates
that the source does not contain; it does not independently verify execution
or completion. The 2013 outcomes are treated as materialized in 2013, although
the SISFOH operation covered 2012--2013. The later outcomes come from the
restricted INEI-assisted Census linkage, not complete community-level Census
tabulations for Peru.

\subsection{Reading the main estimates}

The main tables display the complete registered primary outcome families,
their assignment reduced forms, fuzzy treatment-receipt ratios, robust 95\%
confidence intervals and within-family Holm adjustments. All \PrimaryTests{}
primary common-window fuzzy tests return estimates, and \PrimarySignals{}
have a Holm-adjusted probability below 0.05. The adjustment is within each
registered eight-outcome family, not one global family of all
\PrimaryTests{} tests. This evidence does not establish that reparations had
no effects: many intervals admit substantively important changes in either
direction. The companion community figures show every primary outcome,
standardized by its below-cutoff outcome standard deviation. Original units
remain available in the tables.

The 2013 community outcomes describe rostered population and households,
age composition, social-program participation, wellbeing, unmet basic needs,
labor-force participation and education. Household and person analyses add
granularity but use their own response universes and weights. A community
mean, a household share and a person-level indicator are not interchangeable
estimands. In particular, a larger number of household or person observations
does not supply a corresponding increase in independent treatment-assignment
units. The appendix reports effective observations, represented communities,
first stages and outcome definitions rather than relying on nominal national
sample sizes to convey precision.

The same principle applies to the 2017 community, household and individual
families. Their point estimates should be read jointly with uncertainty,
linkage and completeness. The approved common window
\(h=0.0075\), with bias window \(b=0.0135\), holds the design neighborhood
fixed. It does not make missing-outcome patterns identical or establish that
this window is optimal for every outcome. Registered bandwidth sensitivities
are supporting diagnostics, not alternative windows selected for favorable
findings. Logarithmic outcomes remain in log points: exponentiating a fuzzy
ratio is not automatically a complier-average percentage effect. The GDP
source's community allocation and monetary-unit limitations also remain;
it is not an independent measure of local income inequality or local growth.

\subsection{Observed migration and the population it represents}

The primary individual migration outcome compares valid canonical origin and
destination community codes among linked people with known 2017 age of at
least 14. It does not classify unlinked people as nonmigrants. The fuzzy
estimate is \MigrationEffect{} percentage points, with robust 95\% interval
\([\MigrationLower{},\MigrationUpper{}]\), raw
\(p=\MigrationP{}\) and within-family Holm
\(p=\MigrationHolm{}\). Its effective local sample contains
\MigrationN{} people in \MigrationCCPP{} victim-registry communities.
The matching local-IV Kleibergen--Paap first-stage statistic is
\(F=\MigrationF{}\). Passing the registered \(F>10\) screen is necessary
for the treatment-receipt interpretation used here, but does not settle
continuity, exclusion, monotonicity or selection.

This estimate concerns observed linked adults near the cutoff, conditional
on the identification assumptions. It is not a demonstrated migration effect
for all source residents, nor a precise relative percentage increase. The
other seven individual primary outcomes retain a different complete-case
population. Their effective sample counts therefore differ even though the
geography and nominal score window are common. The migration row is not
silently forced back into the wellbeing-complete population for convenience.

\subsection{Selection is several different processes}

Community entry into the assisted delivery, person linkage, movement
observability, age observability and household outcome completeness are
different selection layers. The supporting exhibits report their parent
frames separately. In the common window the delivery covers 64 of 71 victim
registry communities. It contains 13,010 source people, of whom 9,862 have
observed linkage and movement. All 7,680 linked known-age adults in the
approved local population have observable movement. That conditional 100\%
rate does not establish complete linkage or movement observability for all
source adults: 2017 age is unknown for 3,148 source people, and uncovered
communities are absent from the delivery. Counts are unweighted; the exported
rates give each represented community equal weight within its stated frame.

Household member linkage likewise differs from entry into the full primary
household analysis sample. Of 4,051 source households locally, 3,532 have at
least one linked member, whereas 2,935 have all eight primary household
outcomes. The formal discontinuity in complete-household analysis
availability is \HouseholdSelectionEffect{} percentage points, with robust
95\% interval \([\HouseholdSelectionLower{},\HouseholdSelectionUpper{}]\)
and raw \(p=\HouseholdSelectionP{}\). This is a selection diagnostic,
not a substantive treatment effect surviving a primary outcome-family
correction. Its negative discontinuity makes the household results
selection-sensitive. A different, nonsignificant member-linkage diagnostic
cannot be used to declare linkage ignorable.

The baseline-only propensity exercises remain feasibility diagnostics. They
use score-side terms and registered pre-treatment population, wellbeing and
GDP measures. Positive fitted probabilities do not demonstrate missing at
random, actual population overlap, or transportability. Effective weighted
observation counts are not independent-cluster counts. No propensity-weighted
RD replaces the primary analysis. No linked-adult movement logit is fit when
all eligible responses are observed in its model frame.

The finite-source movement intervals use only known zero-to-one support and
retain missing observations in their denominators. They concern all-age
source people or the fraction of \emph{all source members} of a household,
not the primary eligible-adult population or the primary household fraction
among observed members. They are descriptive support intervals, not sampling
confidence intervals or causal RD bounds. Lee/Dong-type trimming would
require additional population and selection assumptions; the sign of an
availability jump alone does not establish individual monotonic selection.
The current evidence does not warrant dividing these intervals by a first
stage or extrapolating them to missing communities.

\subsection{Heterogeneity: report identification gates before effect contrasts}

The primary treatment-receipt heterogeneity exercise uses local-IV
interactions, not subgroup estimates selected for attractive first stages.
Its instrument set must identify both treatment and the treatment--moderator
interaction. Population, district-capital status and other registered
baseline moderators are substantively motivated, but motivation cannot
replace adequate support or conditional instrument strength. The appendix
retains every registered moderator's support and gate result, and summarizes
all six families without substituting a single first-stage statistic for
the two-endogenous-variable diagnostics.

Most common-window interaction cells fail one or more registered gates.
Gate-passing estimates do not provide a family-adjusted pattern sufficient
to claim a validated moderator of the program's effects. This does not
demonstrate homogeneous effects: weak instruments and limited local
support restrict what can be learned. The separate \texttt{rdhte} results
concern assignment-effect heterogeneity. They cannot be promoted to the main
fuzzy treatment-receipt heterogeneity estimates merely because their
diagnostics appear more favorable.

The supporting overview records \AssignmentSignals{} secondary assignment
contrasts below their original adjusted thresholds among \AssignmentAdjusted{}
supported, adjusted contrasts across the six wave-by-level displays. These
counts combine neither the correction families nor the estimands into one
global test, and related outcomes at different levels are not independent
replications. They warrant interpretation as assignment-response heterogeneity,
not as identified heterogeneity in the effects of receiving reparations.

Project composition and category-specific receipt jumps describe
implementation, not randomized project-type contrasts. Project choice is
post-assignment and endogenous, and its characteristics are observed for
funded communities. Treated-only comparisons do not establish that productive
projects caused migration. The excluded financing-dose estimates remain
excluded: a strong first stage does not establish the additional exclusion
restriction for financing per capita. Population moderation and descriptive
project composition can inform future hypotheses without relabeling these
unidentified contrasts as causal evidence.

\subsection{Mechanism evidence is not an identified indirect effect}

The mechanism exhibits distinguish candidate intermediate-outcome RD ratios
from descriptive associations. They retain the registered multiplicity
families and flag unavailable outcome-sample instrument-strength diagnostics.
Contemporaneous 2017 employment, internet access and observed migration do
not establish mediator-before-outcome ordering. Their associations can also
reflect selection, reverse ordering or treatment-induced common causes.
They therefore do not identify an average causal mediation effect or a
proportion mediated.

The separate associations with 2013 community predictors may help organize
substantive hypotheses, but those predictors are themselves post-treatment
for communities financed earlier. Temporal precedence over a 2017 outcome
alone does not make them valid instruments or ordinary pre-treatment
controls. These results preserve the motivation of the earlier mediation
work while placing its product-of-coefficients interpretation on hold.
The current evidence cannot establish either a liquidity-driven mobility
mechanism or a place-based anchoring mechanism conclusively.

\subsection{What the study can and cannot presently support}

The review documents conditional local estimates with explicit population
boundaries, not an unqualified national impact evaluation. Official program
documentation supports joint A/B priority together with administrative and
geographic conditions, rather than a strict national sequence exhausting A
before B. Dated registration, score and historical eligibility information
is not reconstructed from the later supplied workbook. The owner's decision
to proceed with the CMAN year assumption closes that acquisition prerequisite,
not the associated interpretation limitations.

Fuzzy receipt ratios require assumptions beyond a discontinuity in recorded
treatment. These include credible local continuity, exclusion and
monotonicity; the modern RD treatment is explicit about the additional
conditions for complier interpretation.\footnote{Cattaneo, Idrobo and
Titiunik, \emph{A Practical Introduction to Regression Discontinuity Designs:
Extensions}, author version dated 25 March 2024, Section 3.2, PDF pp.\ 43--45;
\url{https://doi.org/10.1017/9781009441896}.}
Score rounding, limited local support, selection and the historical search
must accompany any eventual causal interpretation. Mass-point adjustments
do not create missing support, and conventional diagnostics cannot prove
all identifying assumptions.

The next publication decision is therefore about the permitted claims and
artifact-level disclosure, not a search for a statistically favorable
replacement design. Every exhibit in this packet remains an internal review
derivative. Neither successful computation nor readable typesetting grants
publication approval, approves canonical main-text roles, or authorizes
replacement of a live Overleaf input.
